% GENERATED FILE. Edit canonical lessons, then run npm run generate:books.
% canonical-source-hash: fnv1a64:35eb1befcfa03672
% canonical-lessons: ML-C07-numbers-1-5, ML-C07-numbers-6-10

\chapter{Numbers One to Ten}
\label{ch:ml-numbers-one-to-ten}

\begin{chapteropening}
\textbf{By the end of this chapter:} \emph{I can count from one to ten in Malayalam.}

Count āṟu to pattu, keep the rare ḻ of ēḻu, and take ompatu apart as 'one short of ten'.
\end{chapteropening}

\section[\ml{ഒന്ന്} \ml{രണ്ട്} \ml{മൂന്ന്} \ml{നാല്} \ml{അഞ്ച്}]{\ml{ഒന്ന്}, \ml{രണ്ട്}, \ml{മൂന്ന്}, \ml{നാല്}, \ml{അഞ്ച്} — one to five}

\label{lesson:ML-C07-numbers-1-5}



\begin{quote}
Of all the Dravidian cousins, Malayalam is Tamil's closest — they were one language a thousand years ago — and its numbers show it.
\end{quote}

\subsection*{What to know first}
\noindent
\begin{tabularx}{\linewidth}{@{}>{\raggedright\arraybackslash}X>{\raggedright\arraybackslash}X>{\raggedright\arraybackslash}X>{\raggedright\arraybackslash}X@{}}
\toprule
\textbf{} & \textbf{numeral} & \textbf{word} & \textbf{said} \\
\midrule
1 & \textbf{\ml{൧}} & \textbf{\ml{ഒന്ന്}} & \emph{onnu} \\
2 & \textbf{\ml{൨}} & \textbf{\ml{രണ്ട്}} & \emph{raṇṭu} \\
3 & \textbf{\ml{൩}} & \textbf{\ml{മൂന്ന്}} & \emph{mūnnu} \\
4 & \textbf{\ml{൪}} & \textbf{\ml{നാല്}} & \emph{nālu} \\
5 & \textbf{\ml{൫}} & \textbf{\ml{അഞ്ച്}} & \emph{añcu} \\
\bottomrule
\end{tabularx}

The symbols \textbf{\ml{൧} \ml{൨} \ml{൩} \ml{൪} \ml{൫}} are Malayalam's own digits.

\begin{cousinweb}
Where Telugu breaks the family with \emph{okaṭi}, Malayalam stays loyal: its \textbf{one} is \emph{onnu}, the direct cousin of Tamil \emph{oṉṟu} (both from the old \emph{oṉ-}). Laid beside Tamil, the whole row is barely changed:

\noindent
\begin{tabularx}{\linewidth}{@{}>{\raggedright\arraybackslash}X>{\raggedright\arraybackslash}X>{\raggedright\arraybackslash}X@{}}
\toprule
\textbf{} & \textbf{Tamil} & \textbf{Malayalam} \\
\midrule
1 & \emph{oṉṟu} & \textbf{onnu} \\
2 & \emph{iraṇṭu} & \textbf{raṇṭu} \\
3 & \emph{mūṉṟu} & \textbf{mūnnu} \\
4 & \emph{nāṉku} & \textbf{nālu} \\
\bottomrule
\end{tabularx}

The one that moved most is \textbf{five}: Tamil \emph{aintu} softened its \emph{nt} to a palatal, giving Malayalam \textbf{añcu} — the same shift English feels between \emph{nature} said stiffly and said naturally.
\end{cousinweb}

\begin{cousinweb}
Every word here ends in \textbf{ chandrakkala} — the little hook \textbf{\ml{്} } on the last letter (\emph{\ml{ഒന്ന്}}, \emph{\ml{നാല്}}) — which trims the final vowel down to a half-heard \emph{ŭ}. It's why these read \emph{onn'}, \emph{nāl'} rather than a full \emph{onnu}, \emph{nālu}. And the doubled letters (\emph{onnu}, \emph{mūnnu}) are real \textbf{held} consonants — linger on the \emph{nn}.
\end{cousinweb}

\subsection*{Your turn}
Say these aloud:

\begin{itemize}
\raggedright
  \item \textquotedblleft{}onnu, raṇṭu, mūnnu, nālu, añcu\textquotedblright{}
  \item the near-identity — \textquotedblleft{}oṉṟu $\leftrightarrow$ onnu\textquotedblright{}, \textquotedblleft{}iraṇṭu $\leftrightarrow$ raṇṭu\textquotedblright{}
  \item the one that shifted — \textquotedblleft{}aintu $\to$ añcu\textquotedblright{}
\end{itemize}

\begin{tcolorbox}[breakable,colback=teal!4,colframe=teal!35!black,title={Before you move on}]
Count to five in Malayalam. (\emph{Onnu, raṇṭu, mūnnu, nālu, añcu}.) Does Malayalam's \textbf{one} follow Tamil or break away like Telugu? (\textbf{Follows Tamil} — \emph{onnu} $\leftrightarrow$ \emph{oṉṟu}.) Which number changed most from Tamil, and how? (\textbf{Five} — \emph{aintu}'s \emph{nt} palatalised to \emph{añcu}.) What does the \textbf{chandrakkala} \ml{്} do to a word's last sound? (Trims its vowel to a half-heard \emph{ŭ}.) Next: six to ten.
\end{tcolorbox}
\section[\ml{ആറ്} \ml{ഏഴ്} \ml{എട്ട്} \ml{ഒമ്പത്} \ml{പത്ത്}]{\ml{ആറ്}, \ml{ഏഴ്}, \ml{എട്ട്}, \ml{ഒമ്പത്}, \ml{പത്ത്} — six to ten}

\label{lesson:ML-C07-numbers-6-10}



\begin{quote}
Malayalam holds onto a Tamil sound the other cousins let go — and you'll hear it in seven.
\end{quote}

\subsection*{What to know first}
\noindent
\begin{tabularx}{\linewidth}{@{}>{\raggedright\arraybackslash}X>{\raggedright\arraybackslash}X>{\raggedright\arraybackslash}X>{\raggedright\arraybackslash}X@{}}
\toprule
\textbf{} & \textbf{numeral} & \textbf{word} & \textbf{said} \\
\midrule
6 & \textbf{\ml{൬}} & \textbf{\ml{ആറ്}} & \emph{āṟu} \\
7 & \textbf{\ml{൭}} & \textbf{\ml{ഏഴ്}} & \emph{ēḻu} \\
8 & \textbf{\ml{൮}} & \textbf{\ml{എട്ട്}} & \emph{eṭṭu} \\
9 & \textbf{\ml{൯}} & \textbf{\ml{ഒമ്പത്}} & \emph{ompatu} \\
10 & \textbf{\ml{൰}} & \textbf{\ml{പത്ത്}} & \emph{pattu} \\
\bottomrule
\end{tabularx}

\begin{cousinweb}
Watch \textbf{seven} travel across the three sister scripts:

\noindent
\begin{tabularx}{\linewidth}{@{}>{\raggedright\arraybackslash}X>{\raggedright\arraybackslash}X>{\raggedright\arraybackslash}X@{}}
\toprule
\textbf{} & \textbf{word} & \textbf{its 7-letter} \\
\midrule
Tamil & \emph{ēḻu} & \textbf{\ta{ழ}} (ḻ) \\
\textbf{Malayalam} & \emph{\textbf{ēḻu}} & \textbf{\ml{ഴ}} (ḻ) — kept \\
Kannada & \emph{ēḷu} & ḷ — hardened \\
Telugu & \emph{ēḍu} & ḍ — hardened further \\
\bottomrule
\end{tabularx}

Malayalam, alone with Tamil, \textbf{keeps the ḻ} — the rare retroflex sound that sits in the name \emph{Tamiḻ} itself. (The name \emph{Malayāḷam}, by contrast, carries the \textbf{hardened} cousin \textbf{ḷ} — \ml{ള} — the very sound Kannada's \emph{ēḷu} shows; the two letters \ml{ഴ} \emph{ḻ} and \ml{ള} \emph{ḷ} are worth keeping apart.) Six (\emph{āṟu}), eight (\emph{eṭṭu}) and ten (\emph{pattu}) are likewise almost letter-for-letter Tamil: this is the closest of all the cousin-columns.
\end{cousinweb}

\begin{cousinweb}
\textbf{Nine} comes, exactly as in Tamil, from the old \emph{oṉ} (\textquotedblleft{}one\textquotedblright{}) + \emph{patu} (\textquotedblleft{}ten\textquotedblright{}) — one short of ten. But before the \emph{p}, the \emph{n} shifts to an \textbf{m}: so the underlying \emph{oṉpatu} is both \textbf{written \ml{ഒമ്പത്}} and \textbf{said \emph{ompatu}} (say \textquotedblleft{}onpatu\textquotedblright{} fast and you'll do it yourself). Kannada does the same in \emph{ombattu}.
\end{cousinweb}

\subsection*{Your turn}
Say these aloud:

\begin{itemize}
\raggedright
  \item \textquotedblleft{}āṟu, ēḻu, eṭṭu, onpatu, pattu\textquotedblright{}
  \item seven's kept sound — \textquotedblleft{}ēḻu\textquotedblright{}, the ḻ of \emph{Tamiḻ} (not the ḷ of \emph{Malayāḷam})
  \item nine taken apart — \textquotedblleft{}oṉ + patu $\to$ ompatu\textquotedblright{}, one short of ten
\end{itemize}

\begin{tcolorbox}[breakable,colback=teal!4,colframe=teal!35!black,title={Before you move on}]
Count six to ten in Malayalam. (\emph{Āṟu, ēḻu, eṭṭu, onpatu, pattu}.) Which sister languages keep Tamil's rare \textbf{ḻ} in \emph{seven}, and which hardened it? (\textbf{Malayalam keeps it}; Kannada $\to$ \emph{ḷ}, Telugu $\to$ \emph{ḍ}.) How is \emph{ompatu} built, and why the \textbf{m}? (\emph{oṉ} \textquotedblleft{}one\textquotedblright{} + \emph{patu} \textquotedblleft{}ten\textquotedblright{}; the \emph{n} becomes \emph{m} before \emph{p}.) Which of these numbers is furthest from Tamil? (None much — this is the closest cousin-set.)
\end{tcolorbox}
